% !TeX root = ../main.tex
% ---------------------------------------------------------------------------
% Ground Improvement and Special Construction Methods
% Scope extension: not derived from the RSE3010 weekly material.
% Key sources: Kirsch & Bell, Ground Improvement; Harris et al. on
% compensation grouting; Andersland & Ladanyi on frozen ground; EN 12715/12716.
% ---------------------------------------------------------------------------
\chapter{Ground Improvement and Special Construction Methods}
\label{ch:ground-improvement}


\section{Objectives and technique selection}
\label{sec:objectives-and-technique-selection}

\subsection{Objectives of ground treatment}
\label{sec:improvement-objectives}
\scopenote{Four distinct aims --- strength increase, stiffness increase,
permeability reduction and movement compensation --- and the frequent error of
selecting a technique before deciding which of them is required.}

\subsection{Applicability by ground type}
\label{sec:improvement-applicability}
\scopenote{The classical grouting applicability chart in grading terms, and
the analogous selection envelopes for jet grouting, freezing and mixing.}

\section{Grouting methods}
\label{sec:grouting-methods}

\subsection{Permeation grouting}
\label{sec:permeation-grouting}
\scopenote{Cement, microfine and chemical grouts; groutability criteria;
sleeve pipe (tube-a-manchette) methods; injection pressure and volume control.}

\subsection{Rock grouting and pre-excavation grouting}
\label{sec:rock-grouting}
\scopenote{Systematic probe-and-grout ahead of the face, grout curtain design,
Lugeon-based acceptance criteria, and its role in inflow control
(Chapter~\ref{ch:hydrogeology}).}

\subsection{Jet grouting}
\label{sec:jet-grouting}
\scopenote{Single, double and triple fluid systems; column geometry
prediction; energy and spoil return; verification coring; and the risks of
uncontrolled heave and spoil blockage.}

\section{Displacement and compensation grouting}
\label{sec:displacement-and-compensation-grouting}

\subsection{Compensation grouting}
\label{sec:compensation-grouting}
\scopenote{Grouting between the tunnel and the structure to intercept
settlement; observational control with real-time monitoring; conditioning,
concurrent and corrective phases; and the limits of what it can protect.}

\subsection{Compaction and fracture grouting}
\label{sec:compaction-grouting}
\scopenote{Displacement grouting for densification and for controlled heave;
distinguishing intended fracture from unintended hydrofracture.}

\section{Freezing and soil mixing}
\label{sec:freezing-and-soil-mixing}

\subsection{Ground freezing}
\label{sec:ground-freezing}
\scopenote{Brine and liquid nitrogen systems; freeze wall design for strength
and cut-off; formation and maintenance time; frost heave and thaw settlement;
energy cost and monitoring.}

\subsection{Deep soil mixing and cutter soil mixing}
\label{sec:soil-mixing}
\scopenote{Panel and column mixing for cut-off walls and block treatment,
strength and uniformity expectations, and quality control.}

\section{Pre-support and face treatment}
\label{sec:pre-support-and-face-treatment}

\subsection{Pre-support ahead of the face}
\label{sec:pre-support}
\scopenote{Pipe roofs and canopy tubes, forepoling spiles, and the
umbrella-arch concept for low-cover and portal work.}

\subsection{Face reinforcement}
\label{sec:face-reinforcement}
\scopenote{Fibreglass dowels in the advance core, the ADECO-RS rationale from
Chapter~\ref{ch:construction}, installation cycle and its effect on advance
rate.}

\section{Dewatering and verification}
\label{sec:dewatering-and-verification}

\subsection{Dewatering systems}
\label{sec:dewatering}
\scopenote{Wellpoints, deep wells, ejectors and sumps; selection by
permeability and drawdown; recharge wells to limit off-site drawdown.}

\subsection{Verification of treated ground}
\label{sec:improvement-verification}
\scopenote{Coring, permeability testing, geophysics and trial panels;
specifying acceptance criteria that can actually be measured before the
treatment is relied upon.}
